\section{Moteur essence (Beau de Rochas) : cycle, rendement, puissance (diapos 133-146)}

\subsection*{Le cycle en 4 temps -- ce qui se passe physiquement}
\begin{demotable}
\textbf{1. Admission} (A$\to$B, isobare) &
La soupape d'admission s'ouvre, le piston descend du PMH au PMB en aspirant un mélange air+carburant, à
pression constante ($p_{atm}$). \\
\textbf{2. Compression} (B$\to$C, \textbf{adiabatique}) &
Soupapes fermées, compression rapide : le système de refroidissement n'a pas le temps d'évacuer la
chaleur, donc $Q\approx0$. \\
\textbf{3. Explosion + détente} (C$\to$D isochore, puis D$\to$E adiabatique) &
La bougie déclenche l'explosion au point C : augmentation \textbf{quasi instantanée} de pression à volume
constant (C-D), puis détente adiabatique (D-E) qui pousse le piston -- c'est le \textbf{temps moteur}. \\
\textbf{4. Échappement} (E$\to$B isochore, puis B$\to$A isobare) &
La soupape d'échappement s'ouvre : chute de pression à volume constant (E-B) puis expulsion des gaz brûlés
lors de la remontée du piston (B-A), à pression constante. \\
\end{demotable}

\begin{center}
\begin{tikzpicture}[scale=1, >=Stealth]
\draw[->] (0,0) -- (6,0) node[right] {$V$};
\draw[->] (0,0) -- (0,5.6) node[above] {$p$};
\coordinate (A) at (1,1); \coordinate (B) at (4.5,1);
\coordinate (C) at (1,3.2); \coordinate (D) at (1,5.2);
\coordinate (E) at (4.5,2.4);
\draw[thick,->] (A) -- (B) node[midway,below]{isobare};
\draw[thick,->] (B) to[out=110,in=-40] (C) node[midway,left,xshift=-3pt]{adiabatique};
\draw[thick,->] (C) -- (D) node[midway,left]{isochore};
\draw[thick,->] (D) to[out=-20,in=140] (E) node[midway,above,xshift=5pt]{adiabatique};
\draw[thick,->] (E) -- (B) node[midway,right]{isochore};
\draw[thick,dashed,->] (B) -- (A) node[midway,below]{isobare};
\foreach \p/\lab/\pos in {A/A/below left,B/B/below right,C/C/above left,D/D/above left,E/E/right}
  {\fill (\p) circle (1.3pt) node[\pos] {\lab};}
\end{tikzpicture}
\end{center}

\begin{retenir}
En réalité, seul le \textbf{cycle utile B-C-D-E-B} compte pour le rendement : l'admission (A-B) et
l'échappement (B-A) sont sur la \textbf{même isobare}, parcourue en sens opposé, donc leur travail s'annule
exactement sur un cycle complet.
\end{retenir}

\subsection*{Démonstration du rendement théorique}
\begin{demotable}
$\eta=1+\dfrac{Q_{EB}}{Q_{CD}}=1+\dfrac{nC_v(T_B-T_E)}{nC_v(T_D-T_C)}$ &
Même logique que Carnot : rendement $=1+\frac{\text{chaleur perdue}}{\text{chaleur reçue}}$. Ici les deux
échanges de chaleur sont isochores, donc $Q=nC_v\Delta T$. \\
Adiabatiques B-C et D-E : $T_BV_B^{\gamma-1}=T_CV_C^{\gamma-1}$ et $T_EV_E^{\gamma-1}=T_DV_D^{\gamma-1}$, avec
$V_E=V_B$ et $V_D=V_C$ (isochores) &
On applique la formule de l'adiabatique aux deux détentes/compressions. Comme la combustion (C-D) et
l'échappement (E-B) sont à volume constant, $V_D=V_C$ et $V_E=V_B$. \\
$\Rightarrow \dfrac{T_B}{T_E}=\dfrac{T_C}{T_D}$ &
Les volumes s'éliminent complètement des deux relations adiabatiques (ils sont identiques deux à deux). \\
$\eta=1+\dfrac{T_B-T_E}{T_D-T_C}=1-\dfrac{T_B-T_E}{T_C-T_D}=\boxed{1-\dfrac{T_B}{T_C}}$ &
\textbf{Astuce} : si $T_B/T_E=T_C/T_D$, alors $(T_B-T_E)/(T_D-T_C)=-T_B/T_C$ (propriété des proportions).
On obtient un rapport de températures tout simple. \\
$\dfrac{T_B}{T_C}=\left(\dfrac{V_C}{V_B}\right)^{\gamma-1}=\left(\dfrac{V_B}{V_C}\right)^{1-\gamma}$ &
On réutilise la relation adiabatique B-C directement pour exprimer ce rapport en fonction des volumes. \\
On pose $\varepsilon=\dfrac{V_B}{V_A}=\dfrac{V_B}{V_C}$ (rapport volumétrique, $8<\varepsilon<10$ pour
l'essence) &
$\varepsilon$ mesure à quel point le mélange est comprimé avant l'explosion. \\
$\boxed{\eta_{\text{théorique}}=1-\varepsilon^{1-\gamma}}$ &
Résultat final : le rendement théorique ne dépend \textbf{que} de $\varepsilon$ et $\gamma$ -- jamais du
carburant, de la vitesse, ou de la puissance ! \\
\end{demotable}

\subsection*{Comment utiliser le rendement, et passer à la puissance}
\begin{demotable}
$\eta_{\text{théo}}=\dfrac{|W_{\text{théo}}|}{Q_{CD}}$, avec $Q_{CD}=m_{carburant}\cdot PCI$ &
Une fois $\eta_{\text{théo}}$ connu (formule ci-dessus), on multiplie par la chaleur \emph{réellement} apportée par
la combustion pour trouver le travail théorique du cycle. \\
$\eta_f=\dfrac{|W_{\text{indiqué}}|}{|W_{\text{théo}}|}$, puis $\eta_{\text{méca}}=\dfrac{|W_{disponible}|}{|W_{\text{indiqué}}|}$ &
On enchaîne les pertes réelles : $\eta_f$ (le cycle réel suit mal les courbes théoriques : combustion pas
instantanée, pertes de charge...), puis $\eta_{\text{méca}}$ (frottements mécaniques). \\
$\eta_{eff}=\eta_{\text{théo}}\cdot\eta_f\cdot\eta_{\text{méca}}=\dfrac{|W_{disponible}|}{Q_{CD}}$ &
Le rendement effectif global est le produit des trois -- c'est lui qui relie directement le carburant brûlé
au travail vraiment récupérable sur l'arbre. \\
$P_{th}=\dfrac{N}{2\cdot60}\cdot W_{th}\cdot z$, \ $P_{eff}=\dfrac{N}{2\cdot60}\cdot W_{dispo}\cdot z$ &
Pour passer du travail (par cylindre, par cycle) à une \textbf{puissance} : on multiplie par le nombre de
cylindres $z$ et par le nombre de cycles par seconde. En 4 temps, \textbf{un cycle = 2 tours}, d'où la
division par $2\cdot60$ ($N$ en tr/min). \\
\end{demotable}

\newpage
\section{PCS et PCI : différence et utilité (diapos 149-150)}

\begin{demotable}
Le \textbf{pouvoir calorifique} est la quantité de chaleur dégagée par la combustion totale de l'unité de
combustible &
C'est «l'énergie contenue» dans 1 kg (ou 1 mole) de carburant, libérée quand il brûle complètement. \\
\textbf{PCS} (pouvoir calorifique \textbf{supérieur}) : quantité de chaleur \textbf{si l'eau créée par la
combustion est sous forme liquide} &
La combustion d'un hydrocarbure produit toujours de l'eau ($C_xH_y+...\to xCO_2+\frac{y}{2}H_2O$). Si on
récupère cette eau condensée (liquide), on récupère aussi sa chaleur de condensation en plus. \\
\textbf{PCI} (pouvoir calorifique \textbf{inférieur}) : quantité de chaleur \textbf{si l'eau créée est
vaporisée} &
Dans un moteur réel, l'eau part sous forme de vapeur avec les gaz d'échappement : on «perd» cette
chaleur de vaporisation. \\
$\boxed{PCI<PCS}$ &
Une partie de la chaleur totale sert à vaporiser l'eau au lieu de chauffer utilement le système -- c'est
pour cela que le PCI est plus faible. \\
\end{demotable}

\begin{retenir}
\textbf{Utilité :} dans les moteurs et chaudières classiques, les gaz d'échappement sortent chauds et
l'eau qu'ils contiennent reste vapeur $\Rightarrow$ on utilise toujours le \textbf{PCI} dans les calculs de
$Q_{comb}=m_{carburant}\cdot PCI$ (formulaire). Le PCS n'est utile que pour les chaudières «à condensation» qui récupèrent aussi la chaleur de condensation des fumées.
\end{retenir}

\section{Excès d'air $\lambda$ : trop ou pas assez, et conséquences (diapos 155-158)}

$$\lambda=\dfrac{\text{quantité d'air réellement aspirée par unité de combustible}}
{\text{quantité d'air stœchiométrique par unité de combustible}}$$

\begin{demotable}
$\lambda$ légèrement $<1$ (pas assez d'air) &
Il n'y a pas tout à fait assez d'oxygène pour bruler tout le carburant $\Rightarrow$ production de
\textbf{monoxyde de carbone} (CO), combustion incomplète et polluante. \\
$\lambda$ fortement $<1$ &
Manque d'air encore plus sévère $\Rightarrow$ présence de \textbf{carburant imbrûlé} dans l'échappement
(gaspillage direct d'énergie et de carburant). \\
$\lambda>1$ (excès d'air) : nécessaire en pratique &
Le mélange air-carburant n'est jamais parfaitement homogène dans le cylindre : il faut $\lambda>1$ pour être
sûr qu'il y a assez d'oxygène \textbf{en tout point} du mélange, même là où il y a localement moins d'air. \\
Mais $\lambda$ trop grand &
Plus $\lambda$ est grand, plus on comprime de l'air «en trop» pour rien : c'est une \textbf{perte
d'énergie} (l'air excédentaire est chauffé et comprimé sans participer à la combustion). \\
\end{demotable}

\begin{retenir}
\textbf{Conclusion du cours :} il faut trouver le bon équilibre -- ni trop peu d'air (CO, imbrûlés,
pollution, perte de rendement) ni trop d'air (perte d'énergie à comprimer de l'air inutile).
\end{retenir}

\newpage
\section{Moteur diesel : cycle complet, comparaison avec l'essence (diapos 157-172)}

\begin{demotable}
Rapport de compression plus important : $16<\varepsilon<20$ &
Contrairement à l'essence ($8<\varepsilon<10$), le diesel comprime beaucoup plus fort. \\
Injection du carburant \textbf{après} la compression &
En essence, air+carburant sont admis \textbf{ensemble}. En diesel, on n'admet que de l'\textbf{air}, et le
carburant est injecté seulement à la fin de la compression. \\
\textbf{Auto-inflammation} du carburant (pas de bougie) &
L'air, très comprimé, est déjà tellement chaud que le carburant injecté s'enflamme tout seul au contact --
pas besoin d'étincelle. \\
\end{demotable}

\subsection*{Les 4 temps, comparés à l'essence}
\begin{demotable}
\textbf{1. Admission} (A$\to$B, isobare) &
Identique à l'essence dans le principe, mais on n'aspire que de \textbf{l'air} (pas de carburant). \\
\textbf{2. Compression} (B$\to$C, \textbf{adiabatique}) &
Identique à l'essence (rapide, pas d'échange de chaleur), mais avec un $\varepsilon$ plus grand : la
température finale $T_C$ est donc bien plus élevée. \\
\textbf{3. Injection + combustion} (C$\to$D, \textbf{isobare}, PAS isochore !) &
Le carburant est injecté en fin de compression et brûle \textbf{progressivement} pendant l'injection, ce qui
se fait à \textbf{pression constante} (contrairement à l'explosion quasi instantanée de l'essence, isochore). \\
\textbf{4. Détente} (D$\to$E, adiabatique) puis \textbf{échappement} (E$\to$B isochore, B$\to$A isobare) &
Identiques à l'essence : détente adiabatique (temps moteur), puis chute de pression à volume constant et
expulsion des gaz brûlés à pression constante. \\
\end{demotable}

\begin{center}
\begin{tikzpicture}[scale=1, >=Stealth]
\draw[->] (0,0) -- (6,0) node[right] {$V$};
\draw[->] (0,0) -- (0,5.6) node[above] {$p$};
\coordinate (A) at (1,1); \coordinate (B) at (4.5,1);
\coordinate (C) at (1,4.2); \coordinate (D) at (2.8,4.2);
\coordinate (E) at (4.5,1.9);
\draw[thick,->] (A) -- (B) node[midway,below]{isobare};
\draw[thick,->] (B) to[out=110,in=-90] (C);
\draw[thick,->] (C) -- (D);
\node at (1.9,4.75) {\small isobare (combustion)};
\draw[thick,->] (D) to[out=-60,in=150] (E);
\draw[thick,->] (E) -- (B) node[midway,right]{isochore};
\draw[thick,dashed,->] (B) -- (A) node[midway,below]{isobare};
\foreach \p/\lab/\pos in {A/A/below left,B/B/below right,C/C/left,D/D/right,E/E/right}
  {\fill (\p) circle (1.3pt) node[\pos] {\lab};}
\end{tikzpicture}
\end{center}

\subsection*{Calcul du rendement (diffère de l'essence à cause de la combustion isobare)}
\begin{demotable}
$\eta=1+\dfrac{Q_{EB}}{Q_{CD}}=1+\dfrac{nC_v(T_B-T_E)}{nC_p(T_D-T_C)}$ &
Comme pour l'essence, mais $Q_{CD}$ utilise maintenant $C_p$ car la combustion C-D est \textbf{isobare} (et
non plus $C_v$). \\
On définit $\rho=\dfrac{V_D}{V_C}$ (rapport de combustion, ou d'injection), en plus de
$\varepsilon=\dfrac{V_B}{V_C}$ &
$\rho$ mesure de combien le volume augmente pendant l'injection/combustion à pression constante (nouveau
paramètre, absent chez l'essence où $V_D=V_C$). \\
Adiabatiques B-C et D-E, plus la relation isobare $T_D/T_C=V_D/V_C=\rho$ &
Mêmes outils que pour l'essence, en ajoutant la loi des gaz parfaits appliquée à l'isobare C-D. \\
\end{demotable}
\eqbox{\eta_{\text{théorique}}=1+\dfrac{\varepsilon^{1-\gamma}-\varepsilon^{1-\gamma}\cdot\rho^{\gamma}}
{\gamma(\rho-1)}}
\emph{En mots :} le rendement théorique du diesel dépend de \textbf{deux} paramètres géométriques
($\varepsilon$ et $\rho$) au lieu d'un seul pour l'essence -- logique, puisque la combustion diesel n'est
plus instantanée mais s'étale sur une phase isobare caractérisée par $\rho$. Les rendements de forme,
mécanique et effectif se calculent \textbf{exactement de la même façon} que pour l'essence.

\newpage
\section{Moteur 2 temps : phases, comparaison avec le 4 temps (diapos 177-182)}

\begin{objectif}
\textbf{Pas de nouvelle démonstration ici} : le cycle thermodynamique idéal est le \textbf{même} que celui
de l'essence (compression adiabatique, combustion isochore, détente adiabatique, échappement) -- donc
\textbf{même formule} $\eta_{\text{théorique}}=1-\varepsilon^{1-\gamma}$. Ce qui change, c'est la
\textbf{mécanique}.
\end{objectif}

\begin{demotable}
Pas de soupapes : deux ouvertures (\textbf{lumières}) dégagées ou bouchées par le piston lui-même &
Plus simple mécaniquement : le piston fait à la fois office de piston et de «distributeur» pour
l'admission/échappement. \\
\textbf{1 tour du moteur = 1 cycle complet} (au lieu de 2 tours en 4 temps) &
Tout se passe en une seule remontée-descente du piston : combustion+détente+échappement+admission se
chevauchent en un seul temps, puis compression dans l'autre. \\
Temps 1 : combustion, détente, échappement \textbf{et} admission simultanés &
Le mélange explose, le piston descend et découvre les lumières : les gaz brûlés sortent pendant que le
mélange neuf entre en même temps (balayage). \\
Temps 2 : compression, pendant que le mélange neuf est aspiré dans le carter &
Le piston remonte et comprime le mélange dans le cylindre, tandis qu'en dessous, le carter aspire déjà le
mélange du \textbf{prochain} cycle. \\
\end{demotable}

\begin{center}
\renewcommand{\arraystretch}{1.5}
\begin{tabularx}{\linewidth}{@{}>{\raggedright\arraybackslash}X >{\raggedright\arraybackslash}X@{}}
\toprule
\textbf{Avantages (vs 4 temps)} & \textbf{Inconvénients (vs 4 temps)} \\
\midrule
À cylindrée égale, puissance plus importante & Bruyant \\
Pour un cylindre, couple moteur plus régulier & Une partie de l'essence passe par l'échappement sans être
brûlée $\Rightarrow$ rendement plus faible \\
Organes simplifiés, moins lourd & Réglage des lumières compliqué \\
\bottomrule
\end{tabularx}
\end{center}

\begin{retenir}
\textbf{Puissance :} $P_{eff}=\dfrac{N}{60}\cdot W_{dispo}\cdot z$ -- \textbf{pas de division par 2} car 1
tour = 1 cycle (contrairement au 4 temps où $P=\frac{N}{2\cdot60}W\cdot z$). C'est le piège classique
à l'examen sur cette matière. Usages typiques : deux-roues, tondeuses -- machines compactes et légères.
\end{retenir}
